\section*{Chapter \ref{ch:g11:production-and-costs} --- Production and Costs}
% figdata: figdata/grade-11/02-production-and-costs/out/ (cereal_marked, fertilizer); ledger g11:dkr2008-*, g11:liberty-*; Kestria numbers invented.

\begin{solution}{exo:g11:production-and-costs:1}
\omterm{def:g11:production-and-costs:marginal-product}{Marginal products}: 10, 15, 11, 8 and 4 tables. The second worker adds more than the first; from the third worker on, the \omterm{def:g11:production-and-costs:marginal-product}{marginal product} is diminishing.
\end{solution}

\begin{solution}{exo:g11:production-and-costs:2}
$600/50 = 12$, $400/50 = 8$ and $200/50 = 4$ kilograms of grain per kilogram of fertiliser.
\end{solution}

\begin{solution}{exo:g11:production-and-costs:3}
At 100: AFC $500/100 = 5$, AVC 3, ATC 8. At 500: AFC 1, AVC 3, ATC 4. The \omterm{def:g11:production-and-costs:atc}{average fixed cost} falls as the fixed cost is spread over more units.
\end{solution}

\begin{solution}{exo:g11:production-and-costs:4}
ATC $300/10 = 30$, $400/20 = 20$ and $600/30 = 20$ dalens. \omterm{def:g11:marginal-reasoning:marginal-benefit}{Marginal cost}: $(400 - 300)/10 = 10$ a unit from 10 to 20, and $(600 - 400)/10 = 20$ a unit from 20 to 30.
\end{solution}

\begin{solution}{exo:g11:production-and-costs:5}
Kenya: about 52 kilograms of fertiliser a hectare and 1.7 tonnes of cereals; the Netherlands: about 275 kilograms and 8.2 tonnes. The two countries also differ in
rain, soil, seeds, irrigation, machines and farmers' training, all of which go with both more fertiliser and higher yields; the chart cannot separate them.
\end{solution}

\begin{solution}{exo:g11:production-and-costs:6}
$\mathit{MC} = 4 + 2q$, $\mathit{AVC} = 4 + q$, $\mathit{ATC} = 100/q + 4 + q$. At 5: 14, 9, 29. At 10: 24, 14, 24. At 20: 44, 24, 29.
\end{solution}

\begin{solution}{exo:g11:production-and-costs:7}
$\mathit{ATC}'(q) = -100/q^2 + 1$, zero at $q = 10$ (negative before, positive after). $\mathit{ATC}(10) = 10 + 4 + 10 = 24 = \mathit{MC}(10)$.
\end{solution}

\begin{solution}{exo:g11:production-and-costs:8}
\omterm{def:g11:marginal-reasoning:marginal-benefit}{Marginal cost} does not change, $4 + 2q$. \omterm{def:g11:production-and-costs:atc}{Average total cost} rises by $300/q$ at every output. Its lowest point moves to $q = \sqrt{400} = 20$, where $\mathit{ATC} = 20 + 4
+ 20 = 44 = \mathit{MC}(20)$.
\end{solution}

\begin{solution}{exo:g11:production-and-costs:9}
Per quarter teaspoon: 28.1 points for the first, 19.5 for the second, and $15.5/2 = 7.75$ for each of the last two. The half teaspoon earned the most, 36.0\,\% over the
season. The whole teaspoon raised the yield most, but its last half teaspoon added maize worth less than it cost.
\end{solution}

\begin{solution}{exo:g11:production-and-costs:10}
\num{80000}, \num{64000}, \num{51200} and $100000 \times 0.8^6 \approx \num{26214}$ hours.
\end{solution}

\begin{solution}{exo:g11:production-and-costs:11}
$1 - 406000/1730000 \approx 1 - 0.235 = 76.5\,\%$. With $0.8^n = 0.235$, $n = \ln 0.235/\ln 0.8 \approx 6.5$ doublings: around the 90th ship, since $2^{6.5} \approx 90$.
\end{solution}

\begin{solution}{exo:g11:production-and-costs:12}
A step that added more yield per quarter teaspoon than the step before: for example, if the whole teaspoon had raised yields by much more than twice the rise from the
half. Because the doses were given at random to plots on the same farms, the plots were alike in soil, rain, seed and the farmer's care, so the differences in yield
can be put down to the fertiliser.
\end{solution}

\begin{solution}{exo:g11:production-and-costs:13}
$\mathit{ATC}(q) = F/q + a + bq$, so $\mathit{ATC}'(q) = -F/q^2 + b$, negative below $\sqrt{F/b}$ and positive above: the minimum is at $q^* = \sqrt{F/b}$. There $F/q^* =
\sqrt{Fb}$ and $bq^* = \sqrt{Fb}$, so $\mathit{ATC}(q^*) = a + 2\sqrt{Fb}$. \omterm{def:g11:marginal-reasoning:marginal-benefit}{Marginal cost} is $a + 2bq^* = a + 2\sqrt{Fb}$, the same.
\end{solution}

\begin{solution}{exo:g11:production-and-costs:14}
Long-run average cost: 12 at 10 (small), 10 at 20 (small), 9 at 30 (medium), 8 at 40 (large). It falls over the whole range: \omterm{def:g11:production-and-costs:economies-of-scale}{economies of scale}. The \omterm{def:g11:production-and-costs:mes}{minimum efficient scale}
is at least 40 units; the data stop before long-run average cost stops falling.
\end{solution}

\begin{solution}{exo:g11:production-and-costs:15}
\omterm{def:g11:production-and-costs:economies-of-scale}{Economies of scale} (the firm moved along a falling long-run average cost curve as it grew); more or better machines per worker; cheaper inputs; a lower quality of
product. Learning ties cost to cumulative output, scale to current output, and capital to the stock of machines: data on all three, year by year, with input prices
and a measure of quality, would let them be told apart, as in the study of the Liberty yard.
\end{solution}

\begin{solution}{pb:g11:production-and-costs:1}
Costs in thousands of kests; average costs in kests a metre.

\textbf{1.} Fixed cost 3600; variable cost $20 \times 60 + 0.25 \times 3600 = 1200 + 900 = 2100$.

\textbf{2.} AFC $3600/60 = 60$, AVC $2100/60 = 35$, ATC 95.

\textbf{3.} $C_M'(q) = 20 + 0.5q$; at 60 it is 50.

\textbf{4.} \omterm{def:g11:marginal-reasoning:marginal-benefit}{Marginal cost} (50) is below \omterm{def:g11:production-and-costs:atc}{average total cost} (95): the average is falling.

\textbf{5.} $\mathit{AVC} = 20 + 0.25q$ rises with $q$: it is lowest at the smallest outputs, close to 20, which is also the \omterm{def:g11:marginal-reasoning:marginal-benefit}{marginal cost} at $q = 0$.

\textbf{6.} $\mathit{ATC}'(q) = -3600/q^2 + 0.25 = 0$ gives $q^2 = 14400$, $q = 120$.

\textbf{7.} $30 + 20 + 30 = 80$.

\textbf{8.} $C_M'(120) = 20 + 60 = 80$.

\textbf{9.} At 100: $36 + 20 + 25 = 81$. At 150: $24 + 20 + 37.5 = 81.5$.

\textbf{10.} 80 kests a metre.

\textbf{11.} $q = \sqrt{900/1.44} = 25$; $\mathit{ATC} = 36 + 20 + 36 = 92$.

\textbf{12.} $q = \sqrt{8100/0.16} = 225$; $\mathit{ATC} = 36 + 20 + 36 = 92$.

\textbf{13.} At 40: small $22.5 + 20 + 57.6 = 100.1$, medium $90 + 20 + 10 = 120$, large $202.5 + 20 + 6.4 = 228.9$: the small mill. At 200: small $4.5 + 20 + 288 = 312.5$, medium $18 +
20 + 50 = 88$, large $40.5 + 20 + 32 = 92.5$: the medium mill.

\textbf{14.} $900/q + 1.44q = 3600/q + 0.25q$ gives $1.19q^2 = 2700$, $q \approx 47.6$, where both cost about 107.5.

\textbf{15.} $3600/q + 0.25q = 8100/q + 0.16q$ gives $0.09q^2 = 4500$, $q \approx 223.6$, where both cost about 92.0.

\textbf{16.} It follows the small mill up to 47.6, the medium mill from 47.6 to 223.6 and the large mill beyond. With three plant sizes it is scalloped:
along the small mill it falls to 92 at $q = 25$ and rises to 107.5 at 47.6, then drops onto the medium mill. Overall it falls to 80 at 120
(\omterm{def:g11:production-and-costs:economies-of-scale}{economies of scale}) and rises after (\omterm{def:g11:production-and-costs:economies-of-scale}{diseconomies of scale}), apart from a small dip to 92 at $q = 225$ along the large mill.

\textbf{17.} The \omterm{def:g11:production-and-costs:mes}{minimum efficient scale} is 120 thousand metres a year, and the lowest long-run average cost is 80 kests a metre.

\textbf{18.} $600/120 = 5$ mills.

\textbf{19.} The medium mill (100 lies between 47.6 and 223.6), at 81 kests a metre.

\textbf{20.} The \omterm{def:g11:production-and-costs:mes}{minimum efficient scale} of a Kestrian textile mill is 120 thousand metres of cloth a year, reached with the medium mill, at a minimum average cost of 80 kests a
metre.
\end{solution}
